\paragraph{Result R58: the original arithmetic quotient from a finite circle.}
The complete periodized-source proof (P1)--(P64), together with the
full source calculations (PSA1)--(PSA29) and (PSD1)--(PSD41), retains
$v_h=2\xi/h$, $w_h(t)=|v_h(1/2+it)|^2/(2\pi)$ and
$m_{h,k}=w_h^{*k}$. Its logarithmic coordinate map keeps the density
$e^{kr+\sum z_i}$ and the ordered-coordinate orientation
$(-1)^{k-1}$. The scalar circle sequence is related to the actual
relative Hilbert fibre by Fourier coefficients
$L^{-1/2}a_n\widehat{\mathcal U_kF_h^{\otimes k}}(2\pi n/L)$,
whose squared norm is $\sum_n\nu_n|a_n|^2$ with
$\nu_n=(2\pi/L)m_{h,k}(2\pi n/L)$. The zero mode remains in this sum.

For the original degree-$N$ polynomial Gram $M_N$ and unchanged
remainder map $J_N$, define the actual finite sampled Gram
$M_N(L,J)$ by retaining $|n|\le J$. The explicit original-source
quantities in (FC3)--(FC3a) produce
$\alpha=1-\eta_P-\eta_T>0$ and $\beta=1+\eta_P$, and prove
\[
 \alpha M_N\preceq M_N(L,J)\preceq\beta M_N,\qquad
 \alpha G_N\preceq G_N(L,J)\preceq\beta G_N,
 \quad G_N=(J_NM_N^{-1}J_N^*)^{-1}.
\]
The second assertion follows by taking minima over the same original
remainder fibres. The two minimizing representatives retain their
explicit difference in $\chi\mathcal P_{N-q}$ and its complete
theta primitive. On the common degree $D=2q+2$, all four endpoint
degrees and both generator raises are covered. Their exact error is
\[
 \left|\mathcal B_{h,k}(L,J)-\mathcal B_{h,k}\right|
      \le2q\log(\beta/\alpha).
\]
At fixed positive error budgets this is $O(q_k)$; its ratio to
$q_k\log k$ tends to zero. The required $L,J$ are explicit functions
of the original weighted and derivative Grams. Their growth and the
arithmetic sample enclosures are retained as actual quantities, rather
than replaced by unproved uniform estimates.

\paragraph{Result R59: the joint critical-jet and circle observation.}
The full density proof obtains the original exponential tails from
the gamma ray integral and the alternating Dirichlet series, including
every canceled value $v_h(\rho)=g^{(m)}(\rho)/h^{(m)}(\rho)\ne0$.
It proves density on the required analytic strip, with its
Cauchy--Schwarz half-width retained, and computes the completed circle
quotient exactly. If $b=p_L$ is the squarefree product of the sampled
roots of $\chi=\chi_{h,k}$, then
\[
 \ker\theta_L=(b)/(\chi),\qquad
 G_N(L)\downarrow G_\infty(L),\qquad
 \operatorname{rad}G_\infty(L)=(b)/(\chi).
\]
For $k\ge2$ the complete difference $G_N(L)-G_\infty(L)$ is
positive definite at every finite admitted degree: a nonzero
minimizing polynomial cannot vanish on the infinite complement of
the selected sample atoms. The limiting kernel retains the full
unsampled primary blocks and all nonconstant jets at sampled roots.

The proof (DC1)--(DC30) compares this with the actual weighted
tensor-critical map $C_k$, whose kernel is $(a)/(\chi)$ for
$a=\chi_{c,k}$. Put $d=\gcd(a,b)$ and $l=\operatorname{lcm}(a,b)$.
Using the complete critical inverse jets and the original weighted
circle coordinates gives
\[
 \operatorname{im}(C_k,\theta_L)
   \simeq\operatorname{im}C_k
        \mathop{\times}_{\mathbb C[S]/(d)}\operatorname{im}\theta_L
   \simeq\mathbb C[S]/(l).
\]
The joint kernel is $(l)/(\chi)$ and the common quotient is
$\mathbb C[S]/(d)$. At each original sum root, their quotient
lengths are respectively the maximum and minimum of the critical
cyclic length and the sampled length in $\{0,1\}$; all local unit
cofactors are retained in the kernel inclusions. A factorization
$\theta_L=T C_k$ exists exactly when $b\mid a$, with its explicit
inverse-unit and reduction formula. The reverse factorization exists
exactly when $a\mid b$. At $k=1$, full-order cancellation proves
that every sampled selected root has positive weight and $b\mid a$.
For general tensor degree, the extra circle summand has defining
polynomial $b/d$ and the explicit original polynomial source lift
in (DC29). Its existence criterion asserts no existence of an
off-critical zero.

All four observations transport through the same original period
matrix. For the role $s\in\{a,b,d,l\}$, retain its actual map $T_s$
and multiplier $M_s$. Then
\[
 T_{s,t}=T_s\Pi^{-1},\qquad
 u\partial_tT_{s,t}-M_sT_{s,t}
       =tT_s(e_0)\ell\Pi^{-1}.
\]
Their kernels, intersection, sum, crossed observations and complete
rank-one action defects are calculated in (DC23)--(DC30). Their
Split-Zero lifts retain each support and its receiving supported zero.

\paragraph{Result R60: universal residue endomorphisms and metric variation.}
The full calculation (RCX1)--(RCX40) extends the residue insertion
$R=e_0\ell$ over every commutative coefficient ring $B$ with monic
$\chi$. With the original residue matrix
$\mathcal B_{ij}=\ell(S^{i+j})$, its determinant remains
$(-1)^{q(q-1)/2}$. Every coefficient endomorphism has the exact inverse
expansion
\[
 T=\sum_{i,j}(T\mathcal B^{-1})_{ij}A^iRA^j.
\]
The full proof gives matrix units, their multiplication and base
change. Their original source lifts $r_NTj_E$ form an injective
representation in the corner with identity $r_Nj_E$ and kill the
same original theta boundaries.

For a proper nonzero invariant constituent $I:F\hookrightarrow E$,
retain the residual $r_M=(1-P_M)e_0$ and the residue-dual vector
$w_M=I(I^*MI)^{-1}(\ell I)^*$. They are $M$-orthogonal, with
squared norms $q_M,d_M>0$, and
$\mathfrak c_F(M)=q_Md_M$. Along any actual differentiable
positive-form path,
\[
 \frac{d}{ds}\log\mathfrak c_F(M)
 =\operatorname{Tr}\left[
   \left(\frac{r_Mr_M^*M}{q_M}-\frac{w_Mw_M^*M}{d_M}\right)
                  M^{-1}\dot M\right].
\]
The parenthesized operator has eigenvalues $1,-1,0$ with the
stated multiplicities, yielding the exact spectral-spread bound.
For the original endpoint forms $G_i\succeq G_j$, with extreme
relative eigenvalues $0<\theta\le\eta\le1$, this proves
\[
 \left|\log\frac{\mathfrak c_F(G_j)}{\mathfrak c_F(G_i)}\right|
 \le\log(\eta/\theta)\le-\log\theta
 \le\log\frac{\det G_i}{\det G_j}.
\]
The complete neighboring-minor formula
$\mathfrak c_F(M)=v_{J_+}(M)v_{J_-}(M)/v_{J_g}(M)^2$
retains every original frame determinant and both constituent
dimensions. The finite sampling comparison from R58 gives the
dimension-free coefficient interval
$\alpha/\beta\le\mathfrak c_F(G_N(L,J))/\mathfrak c_F(G_N)
\le\beta/\alpha$.

\paragraph{Result R61: sampled control with its exact commutator error.}
The full proof (FC1)--(FC29), including (FC3a), keeps the two actual
forms $G=G_N$, $H=G_N(L,J)$, the original
$T_t=A+tR-k/2$, and $B=G^{-1}(H-G)$. It first proves
\[
 (I+B)X_H=X_G+T_t^{\dagger_G}B+BT_t,
 \qquad X_G=T_t^{\dagger_G}+T_t,
 \quad X_H=H^{-1}(T_t^*H+HT_t).
\]
The coordinating Zeta calculation sharpens the transfer through
the explicit isometry $\mathcal S=(I+B)^{1/2}:(E,H)\to(E,G)$.
Set $Z=[\mathcal S,T_t]\mathcal S^{-1}$. Its complete Hermitian
error and the convergent Sylvester integral give
\[
 \left|\|X_H\|_H-\|X_G\|_G\right|
 \le\|Z+Z^{\dagger_G}\|_G
 \le\alpha^{-1}\|[B,A]+t[B,R]\|_G,
 \qquad [B,R]v=(Be_0)\ell(v)-e_0\ell(Bv).
\]
Both original metrics, the original unit and the ordered residue
terms remain. Commuting discrepancies give exactly equal control
radii. For the actual Fourier-projected representative
$\mathscr R$ and its theta boundary $\mathscr B$, the full identities
are
\[
 \begin{gathered}
 W_H(0)=-(\mathscr R^*\mathscr B+\mathscr B^*\mathscr R),\\
 (D_L^2-kD_L)\mathscr R-\mathscr R(A^2-kA)
       =(D_L-k)\mathscr B+\mathscr BA.
 \end{gathered}
\]
The resulting original-metric control radius gives both inequalities
of the full Laplacian numerical-range parabola (FC18), including the
zero-radius case. These statements preserve the actual arithmetic
commutator and boundary that a numerical endpoint estimate must
evaluate.

\paragraph{Result R62: quotient curvature and the complete fourth derivative.}
Retain the exact original coefficient sequence
$0\to F\xrightarrow I E\to E/F\to0$ and a constant quotient lift
$J_F$. Its least-period-norm lift and quotient Gram $K(t)$ are the
full Schur formulas (FC21). With $C=[I,J_F]$ they satisfy
\[
 \det(I^*\Pi^*\Pi I)\det K
      =|\det C|^2|\det\Pi|^2,
 \qquad
 \partial_{\bar t}\partial_t\log\det K
      =-\frac{|t|^2}{|u|^2}\mathfrak c_F(\Pi^*\Pi).
\]
The positive constituent curvature is therefore compensated in this
same coefficient quotient, with both curvatures vanishing at zero
and with opposite nonzero leading coefficients. The full proof
includes quotient-lift independence, basis covariance and the
literal comparison with the original theta quotient metric.

The residue continuation evaluates the pure jets as well. If
$\psi=\log\det(I^*\Pi(t)^*\Pi(t)I)$ and
$\phi(t)=\log\det(I^*\Pi(0)^*\Pi(t)I)$ near zero, then
\[
 \left.\frac{d^4}{dx^4}\psi(x)\right|_0
   =2\operatorname{Re}\phi^{(4)}(0)
       +6\mathfrak c_F(\Pi(0)^*\Pi(0))/|u|^2.
\]
The full ordered recursion for $\Pi^{(n)}=\Pi Q_n$, its explicit
$Q_4$ including $3R^2/u^2$, and the complete trace-log formula in
(RCX30)--(RCX32) evaluate the remaining pure term. The forced
period equation retains $u\Pi R$, and its original source
primitive retains $tKR$. The exact Pascal translation preserves
these derivatives of order at least two while its scalar factor
retains the first-derivative trace shift $2p\operatorname{Re}a$.
